\documentclass[conference]{IEEEtran}
\IEEEoverridecommandlockouts
\usepackage{cite}
\usepackage{amsmath,amssymb,amsfonts}
\usepackage{algorithmic}
\usepackage{graphicx}
\usepackage{textcomp}
\usepackage{xcolor}

\begin{document}

\title{Quantum SupplyChain Intelligence (QSCI): A Hybrid Quantum-Classical Decision Architecture for NP-Hard Operations\\
\thanks{This work was developed as part of the QSCI Research Initiative.}
}

\author{\IEEEauthorblockN{1\textsuperscript{st} QSCI Team}
\IEEEauthorblockA{\textit{AI & Quantum Lab} \\
\textit{QSCI Enterprise}\\
City, Country \\
email@qsci.com}
}

\maketitle

\begin{abstract}
Supply chain optimization involves several NP-Hard problems including vehicle routing, multi-echelon inventory optimization, and network design. While classical solvers (e.g., Gurobi, OR-Tools) and classical ML models (e.g., TFT, LightGBM) define the current state-of-the-art, Noisy Intermediate-Scale Quantum (NISQ) algorithms offer theoretical scaling advantages. In this paper, we present the Quantum SupplyChain Intelligence (QSCI) platform—a production-grade enterprise architecture that utilizes classical AI as the operational baseline while seamlessly integrating a Quantum Co-Processor abstraction for honest, reproducible benchmarking. We introduce a novel 4-factor recommendation confidence framework, integrate Human-in-the-Loop continuous learning, and evaluate Variational Quantum Regressors (VQR) and Quantum Approximate Optimization Algorithms (QAOA) against their classical counterparts on real-world datasets. Our benchmarking framework explicitly accounts for QRAM data loading costs to ensure fair comparisons.
\end{abstract}

\begin{IEEEkeywords}
Quantum Machine Learning, Supply Chain Optimization, Decision Intelligence, VQR, QAOA
\end{IEEEkeywords}

\section{Introduction}
Global supply chains account for approximately \$19.3T annually, yet decision-making remains heavily reliant on manual processes or opaque ML models. The QSCI platform addresses the "Trust Problem" through Explainable AI (XAI) and Human-in-the-Loop (HITL) feedback, while simultaneously addressing the "Quantum Honesty Problem" by treating quantum hardware as a research co-processor rather than claiming premature quantum advantage.

\section{System Architecture}
The QSCI architecture is bifurcated into a Classical Production Baseline and a Quantum Research Lab. The classical baseline utilizes Temporal Fusion Transformers (TFT) for forecasting and Gurobi/OR-Tools for optimization. The Quantum Lab implements Qiskit Runtime V2 abstractions to execute QAOA, VQE, and VQR on IBM Heron processors.

\section{Benchmarking Methodology}
To ensure fair comparison, our QuantumFairnessFramework penalizes quantum wall-clock time by $O(N)$ to account for QRAM state preparation.

\section{Conclusion}
While classical methods remain dominant for supply chain execution today, the QSCI architecture provides the enterprise scaffolding required to seamlessly transition to quantum dominance when hardware matures.

\end{document}
